% Figure: Gaussian-beam envelope. The beam radius w(z) is a hyperbola with a
% minimum waist w0 at z=0, a near-collimated region within the Rayleigh range
% +-zR, and a far-field cone whose half-angle is the divergence theta. The
% asymptotes show the far-field cone.
\begin{figure}[htbp]
\centering
\begin{tikzpicture}
\begin{axis}[
  width=13cm, height=6.6cm,
  xlabel={axial distance $z$ (units of $z_R$)},
  ylabel={beam radius (units of $w_0$)},
  xmin=-4, xmax=4, ymin=-4.2, ymax=4.2,
  axis lines=middle,
  xtick={-3,-2,-1,0,1,2,3},
  ytick={-3,-2,-1,1,2,3},
  tick label style={font=\scriptsize},
  label style={font=\small},
  samples=200, domain=-4:4,
  legend style={font=\scriptsize, at={(0.02,0.97)}, anchor=north west, draw=none},
]
% w(z)/w0 = sqrt(1 + z^2) with z in units of zR
\addplot[engblue, very thick] {sqrt(1 + x^2)};
\addlegendentry{$+w(z)$}
\addplot[engblue, very thick] {-sqrt(1 + x^2)};
% far-field asymptotes: slope = 1 in these units (w ~ |z| for large z)
\addplot[engorange, dashed, domain=0.6:4] {x};
\addplot[engorange, dashed, domain=-4:-0.6] {-x};
\addplot[engorange, dashed, domain=0.6:4] {-x};
\addplot[engorange, dashed, domain=-4:-0.6] {x};
% Rayleigh range markers
\draw[engred, dotted] (axis cs:1,-4.2) -- (axis cs:1,4.2);
\draw[engred, dotted] (axis cs:-1,-4.2) -- (axis cs:-1,4.2);
% waist arrow
\draw[englime!55!engblack, <->, thick]
  (axis cs:0,-1) -- (axis cs:0,1);
\node[englime!55!engblack, font=\scriptsize, anchor=west]
  at (axis cs:0.08,0.55) {$2w_0$};
\node[engred, font=\scriptsize, anchor=south] at (axis cs:1,3.4)
  {$z=+z_R$};
\node[engorange, font=\scriptsize, anchor=south west] at (axis cs:2.6,2.4)
  {far-field cone $\theta$};
\end{axis}
\end{tikzpicture}
\caption[Gaussian-beam envelope]{The radius $w(z)=w_0\sqrt{1+(z/z_R)^2}$ of a
Gaussian beam against axial distance, scaled to the waist $w_0$ and the
Rayleigh range $z_R=\pi w_0^2/\lambda$. Inside $\pm z_R$ the beam stays within
a factor $\sqrt{2}$ of its minimum radius, the working ``collimated'' length.
Beyond a few $z_R$ the envelope approaches the straight far-field cone of
half-angle $\theta=\lambda/(\pi w_0)$, the divergence that ties a tight waist
to a fast spread. A small waist buys a short collimation length and a wide
far-field cone; a large waist buys the reverse.}
\label{fig:vol04ch11:gaussian-beam}
\end{figure}
